% -*- coding: utf-8 -*-
% Guía de las aplicaciones web para los alumnos (las de ~/Textos/AplicacionesLogica).
% Documento aparte: no forma parte del libro, que remite a las aplicaciones con la macro
% \urlaplicaciones de preambulo_comun.tex. Sustituye (9-11-2026) a la guía de ndprover.org,
% archivada en archivo/guia_verificador_ndprover.tex. Revisarla si cambian las aplicaciones.
\documentclass[a4paper,spanish,11pt]{article}
% -*- coding: utf-8 -*-
% Preámbulo común para la asignatura Lógica (Grado en Filosofía, Universidad de Sevilla)
% Compatible tanto para compilación independiente (article) como en el manual conjunto (book/report)

\usepackage[T1]{fontenc}
\usepackage[utf8]{inputenc}
\usepackage[spanish,es-tabla]{babel}
\usepackage{amsmath,amssymb,amstext,amsthm,array}
\usepackage[top=20mm,bottom=20mm,left=20mm,right=20mm,headsep=5mm,footskip=10mm]{geometry}
\usepackage{enumerate}
\usepackage{tikz}
\usetikzlibrary{trees}

% Paquete Fitch y reglas canónicas de la asignatura
% Altura de línea algo menor que la predeterminada (4.5ex/3.5ex/1.5ex) para que
% las derivaciones sean más compactas y den menos problemas de paginación.
\usepackage[justsep=1em,rules=myrules,justformat=myjust,height=3.8ex,topheight=3.2ex,depth=1.3ex]{fitch}
\newcommand{\myjust}[2]{#1 #2}
\newcommand{\myrules}{
 \ndrules
 \def\ci{\by{$\to$I}}
 \def\ce{\by{$\to$E}}
 \def\x{\by{X}}
 \def\ip{\by{IP}}
 \def\ds{\by{DS}}
 \def\mt{\by{MT}}
 \def\lem{\by{LEM}}
 \def\dne{\by{DNE}}
 \def\dem{\by{DeM}}
 \def\cq{\by{CQ}}
 \def\bi{\by{$\leftrightarrow$I}}
 \def\be{\by{$\leftrightarrow$E}}
 \def\idi{\by{$=$I}}
 \def\ide{\by{$=$E}}}

% Enlaces hipertexto
\usepackage[hidelinks,colorlinks=true,linkcolor=blue!60!black,urlcolor=blue!60!black]{hyperref}

% Aplicaciones web de apoyo (GitHub Pages del repositorio fersoler/AplicacionesLogica).
% Si cambia la dirección, basta cambiarla aquí.
\newcommand{\urlaplicaciones}{https://fersoler.github.io/AplicacionesLogica/}
\newcommand{\webaplicaciones}{\expandafter\url\expandafter{\urlaplicaciones}}
\newcommand{\enlaceapp}[2]{\href{\urlaplicaciones #1/}{\emph{#2}}}
\newcommand{\appsintaxis}{\enlaceapp{sintaxis}{Sintaxis de las fórmulas}}
\newcommand{\appdeduccion}{\enlaceapp{deduccion}{Deducción natural}}
\newcommand{\apptablas}{\enlaceapp{tablas}{Tablas de verdad}}
\newcommand{\appkalmar}{\enlaceapp{kalmar}{La prueba de Kalmár}}
\newcommand{\appmundo}{\enlaceapp{mundo}{El mundo de Tarski}}
\newcommand{\apptableaux}{\enlaceapp{tableaux}{Tableaux semánticos}}
% Nota breve tras una explicación: con qué aplicación se practica. Empieza con
% el icono de un portátil (fontawesome5) en el color de los enlaces.
\usepackage{fontawesome5}
\newcommand{\practica}[1]{\par\smallskip\noindent{\small{\color{blue!60!black}\faLaptop}\enspace\textit{Para practicar.} #1\par}\smallskip}

% Teoremas y entornos formales
\ifcsname chapter\endcsname
  \newtheorem{theorem}{Teorema}[chapter]
\else
  \newtheorem{theorem}{Teorema}[section]
\fi
\newtheorem{lemma}[theorem]{Lema}
\newtheorem{definition}[theorem]{Definición}
\newtheorem{definicion}[theorem]{Definición}
\newtheorem{example}[theorem]{Ejemplo}
\newtheorem{ejemplo}[theorem]{Ejemplo}
\newtheorem{corollary}[theorem]{Corolario}
\newtheorem{proposition}[theorem]{Proposición}
\newtheorem{observacion}[theorem]{Observación}
\newtheorem{nota}[theorem]{Nota}
\newtheorem{aplicacion}[theorem]{Aplicación}

\newenvironment{resumenbox}{%
  \begin{quote}\small
}{%
  \end{quote}
}

% Ejercicios: en el manual (clase book) se numeran por tema (Ejercicio 4.12).
% El título entre paréntesis va en redonda, como en los teoremas y ejemplos.
\newcounter{preg}
\ifcsname chapter\endcsname
  \makeatletter\@addtoreset{preg}{chapter}\makeatother
  \renewcommand{\thepreg}{\thechapter.\arabic{preg}}
\fi
\newenvironment{pregunta}[1][]{%
  \par\addvspace{0.4cm}\refstepcounter{preg}\noindent\textbf{Ejercicio \thepreg}\if\relax\detokenize{#1}\relax\textbf{.}\else\ (#1)\textbf{.}\fi\enspace\ignorespaces
}{%
  \par\vspace{0.4cm}%
}

% Soluciones (apéndice de soluciones)
% Dentro de una solución no se permite un salto de página entre el título y lo
% que le sigue (por ejemplo, una derivación en un entorno center), ni que un
% párrafo deje una sola línea al pie o al principio de una página.
\newenvironment{solucion}[1]{%
  \par\medskip
  \csname @beginparpenalty\endcsname=10000\relax
  \clubpenalty=10000\relax\widowpenalty=10000\relax
  \noindent\textbf{Solución del ejercicio~\ref{#1}.}\enspace\ignorespaces
}{%
  \par\medskip%
}

% Derivación dentro de un recuadro de regla (tabla con \hline): deja un poco
% de aire para que la línea vertical de la derivación no toque las líneas
% horizontales de la tabla.
\newcommand{\cajaregla}[1]{\vbox{\vskip 5pt\hbox{#1}\vskip 4pt}}

% Marca para secciones de ampliación (no necesarias para seguir el curso)
\newcommand{\ampliacion}{\texorpdfstring{\,$^{\star}$}{ (*)}}

\newcounter{pregexamen}
\newenvironment{preguntaexamen}[2]{%
  \refstepcounter{pregexamen}%
  \vspace{0.4cm}%
  \noindent\textbf{Pregunta \arabic{pregexamen}} \hfill \textbf{[#1 puntos]}\\[0.15cm]%
  \noindent\textit{#2}\\[0.2cm]%
}{%
  \vspace{0.3cm}%
}

\newtheorem{ejercicio}{Ejercicio}

\newcounter{ejerciciocnt}
\newcommand{\ejer}[1]{\vspace{1.5em}\noindent\stepcounter{ejerciciocnt}\textbf{Ejercicio \theejerciciocnt.} #1\vspace{0.5em}}

\newenvironment{resumentema}{\begin{quote}\small\textbf{Resumen.}\enspace}{\end{quote}\bigskip}

\renewcommand{\arraystretch}{1.25}
\newcommand{\tecla}[1]{\texttt{#1}}

\begin{document}

\begin{center}
  {\Large \textbf{Aplicaciones web para practicar}}\\[0.15cm]
  {\large Lógica, Grado en Filosofía --- Curso 2026/2027}\\[0.1cm]
  {\small \webaplicaciones}
\end{center}

Las aplicaciones funcionan en el navegador, también en el móvil, sin instalar nada ni registrarse. Usan la notación de los apuntes y cada una tiene al final de la página una ayuda desplegable con lo que aquí se resume.

\section{Las aplicaciones}

\begin{center}
\small
\begin{tabular}{p{3.6cm}p{1.2cm}p{10.2cm}}
\hline
\textbf{Aplicación} & \textbf{Tema} & \textbf{Qué se practica} \\
\hline
\appsintaxis & 1 y 4 & Conectiva principal, subfórmulas, árbol de formación y convenciones de precedencia; alcance de los cuantificadores, variables libres y ligadas, y sustitución. Para explorar una fórmula o hacer ejercicios generados al azar. \\
\appdeduccion & 1 y 4 & Derivaciones con las reglas de los apuntes, comprobadas línea a línea. Contiene los ejercicios de deducción natural de los temas 1 y 4, sin soluciones. \\
\apptablas & 2 & Tablas de verdad para rellenar, clasificación de fórmulas, satisfacibilidad, equivalencia, consecuencia con contraejemplos y asignación directa, evaluación en una interpretación y formalización. \\
\appkalmar & 3 & La prueba de completud: la tabla de verdad, la derivación de cada fila según el lema de Kalmár y el colapso de las premisas con LEM, paso a paso. \\
\appmundo & 5 & Un tablero de piezas que funciona como estructura: evaluación de fórmulas con la definición de satisfacción, contramodelos y formalización. \\
\apptableaux & 6 & Tableaux con las reglas $\alpha$, $\beta$, $\gamma$, $\delta$ y $\sigma$ de los apuntes, paso a paso o construidos por ti, con el modelo de cada rama abierta. \\
\hline
\end{tabular}
\end{center}

\noindent \apptablas{} y \apptableaux{} resuelven por sí mismas lo que se les da. Intenta antes el ejercicio a mano, o con la opción «Construirlo yo» de los tableaux, y úsalas para comprobarlo.

\section{Cómo escribir las fórmulas}

Cada símbolo se escribe con su botón, con el de los apuntes o con el teclado. Cada aplicación admite los que necesita: los cuantificadores y la identidad, solo las de primer orden.

\begin{center}
\begin{tabular}{lcl}
\hline
\textbf{Símbolo} & \textbf{Apuntes} & \textbf{Teclado} \\
\hline
negación & $\neg$ & \tecla{\textasciitilde} \\
conjunción & $\land$ & \tecla{\&} \\
disyunción & $\lor$ & \tecla{|} \quad (\textbf{no} \tecla{v}, que es una variable) \\
condicional & $\to$ & \tecla{->} \\
bicondicional & $\leftrightarrow$ & \tecla{<->} \\
contradicción & $\bot$ & \tecla{\#} \\
cuantificador universal & $\forall x$ & \tecla{Ax} \\
cuantificador existencial & $\exists x$ & \tecla{Ex} \\
identidad & $a = b$ & \tecla{a = b} \\
distinto & $a \neq b$ & \tecla{a != b} \\
\hline
\end{tabular}
\end{center}

\begin{itemize}
    \item Las \textbf{variables proposicionales} son letras sueltas: $p$, $q$, $r$\dots
    \item En primer orden, las \textbf{constantes} son las minúsculas de la \tecla{a} a la \tecla{r}, y las \textbf{variables}, de la \tecla{s} a la \tecla{z}. Los apuntes usan $s$ para Sócrates: en las aplicaciones es una variable, así que usa otra constante.
    \item Los \textbf{predicados} se escriben con paréntesis, \tecla{R(x, b)}, o con los argumentos pegados, \tecla{Rxb}; en la forma pegada, \tecla{Ax} y \tecla{Ex} son cuantificadores. En \appmundo{} los predicados son palabras de su diccionario, como \tecla{IzqDe(a, x)}.
    \item Las \textbf{convenciones} son las de los apuntes (sección 1.1.1): la precedencia, de más a menos fuerte, es $\neg$, $\land$, $\lor$, $\to$, $\leftrightarrow$, y las conectivas binarias se agrupan por la derecha. Los cuantificadores, como el negador, afectan solo a la fórmula más breve que les sigue: para que $\forall x$ abarque un condicional hay que escribir $\forall x\, (P(x) \to Q(x))$.
    \item Las aplicaciones muestran lo que escribes con la notación de los apuntes; así ves cómo lo han leído. En \appdeduccion{}, cada fórmula se reescribe al salir de su campo.
\end{itemize}

\section{Deducción natural}

\begin{itemize}
    \item \textbf{El problema.} Elige un ejercicio de la lista o escribe tu propio argumento en «Otro argumento», con una premisa por línea. Algunos ejercicios traen la pista de los apuntes, oculta hasta que la pides, y los de plantilla guiada pueden cargarla.
    \item \textbf{Las líneas.} «Nueva línea» (\tecla{Intro}) añade una línea en el mismo nivel; «Abrir caja» (\tecla{Ctrl+Intro}) abre una caja con una hipótesis; «Cerrar caja» (\tecla{Mayús+Intro}) la cierra y sigue fuera. Al insertar o borrar líneas, las citas se renumeran solas.
    \item \textbf{Las hipótesis} se justifican con «Hip.», «Hip. IP» o «Hip. testigo $c$». En las demás líneas se escribe la regla y las líneas que cita, separadas por comas; una caja se cita de su hipótesis a su última línea, con guion. Mientras escribes una justificación, debajo de la prueba aparecen los botones de las reglas.
\end{itemize}

\begin{center}
\small
\begin{tabular}{lll}
\hline
\textbf{Regla} & \textbf{Con el teclado} & \textbf{Qué se cita y en qué orden} \\
\hline
R & \tecla{R 1} & la línea copiada \\
$\land$I, $\land$E & \tecla{\&I 1, 2}, \tecla{\&E 1} &  \\
$\to$I & \tecla{->I 2-5} & la caja \\
$\to$E & \tecla{->E 1, 2} & primero el condicional, luego el antecedente \\
$\lor$I & \tecla{|I 1} &  \\
$\lor$E & \tecla{|E 1, 2-4, 5-7} & primero la disyunción, luego las dos cajas \\
$\neg$I & \tecla{\textasciitilde I 2-4} & la caja \\
$\neg$E & \tecla{\textasciitilde E 3, 2} & primero la negación, luego la fórmula \\
X & \tecla{X 4} & la línea con $\bot$ \\
IP & \tecla{IP 1-7} & la caja que empieza con la negación \\
$\leftrightarrow$I, $\leftrightarrow$E & \tecla{<->I 3-5, 6-8}, \tecla{<->E 1, 2} & las dos cajas; en $\leftrightarrow$E, primero el bicondicional \\
MT, DS & \tecla{MT 1, 2}, \tecla{DS 1, 2} & primero el condicional; primero la disyunción \\
DNE, LEM, DeM & \tecla{DNE 3}, \tecla{LEM 3-4, 5-6}, \tecla{DeM 2} & en LEM, las dos cajas \\
$\forall$E, $\forall$I & \tecla{AE 1}, \tecla{AI 8} &  \\
$\exists$I, $\exists$E & \tecla{EI 5}, \tecla{EE 2, 3-6} & en $\exists$E, primero el existencial, luego la caja \\
CQ & \tecla{CQ 1} &  \\
=I, =E & \tecla{=I}, \tecla{=E 1, 3} & =I no cita nada; en =E, primero la identidad \\
\hline
\end{tabular}
\end{center}

\begin{itemize}
    \item \textbf{El orden de las citas cuenta}: es el de los esquemas de los apuntes (apéndice C).
    \item \textbf{Comprobación.} Cada línea muestra \checkmark{} o $\times$, con el motivo del error, y debajo se indica si la prueba está completa. Se comprueban las restricciones de $\forall$I y $\exists$E tal como las dan los apuntes, y que todas las líneas sean sentencias.
    \item \textbf{Restricciones de los ejercicios.} La casilla «Solo reglas básicas» rechaza MT, DS, DNE, LEM, DeM y CQ, y quita sus botones. Se marca sola en los ejercicios que lo piden (1.15, 1.22, 4.16 y 4.17). En el 1.18 y el 1.23, que piden derivar dos leyes de De Morgan, no se admite DeM.
    \item \textbf{Sin soluciones.} La aplicación no genera pruebas ni muestra soluciones: lo que hay que aprender es a construir la prueba.
    \item \textbf{Exportar.} La prueba se exporta a \LaTeX{}, con el paquete \texttt{fitch} de los apuntes, o a texto.
    \item \textbf{Guardado.} La prueba se guarda solo en tu navegador y en la dirección de la página: copiarla sirve para guardarla o enviarla. No se envía nada a ningún sitio ni hay cookies. La ayuda de la página tiene un botón para borrar lo guardado.
\end{itemize}

\section{Tableaux semánticos}

\begin{itemize}
    \item \textbf{Qué se comprueba} (sección 6.5): consecuencia, validez, satisfacibilidad de un conjunto o equivalencia. La aplicación pone en la raíz las fórmulas que corresponden, por ejemplo las premisas y la negación de la conclusión.
    \item \textbf{Verlo paso a paso.} El árbol se construye con la prioridad de la sección 6.3.2 ($\alpha$ y $\sigma$, $\delta$, $\gamma$, $\beta$), expandiendo primero la rama abierta de más a la izquierda, y cada paso se explica.
    \item \textbf{Construirlo yo.} Pulsas la fórmula que quieres descomponer y la regla se aplica en todas las ramas abiertas que pasan por ella; en $\gamma$ eliges la constante. Las ramas las cierras tú: pulsa $\otimes$ al final de la rama y señala los dos nodos contradictorios. Si no sigues la prioridad, o sigues descomponiendo en una rama que ya podías cerrar, la aplicación te lo comenta.
    \item \textbf{El modelo.} De cada rama abierta y completa da el modelo, como en la sección 6.4, y comprueba que satisface la raíz.
    \item \textbf{Ramas infinitas} (sección 6.9). Al llegar al límite de constantes por rama, la rama queda incompleta. La opción «Regla $\delta$ de Boolos» permite encontrar modelos finitos.
    \item No admite identidad ni funciones, como el tema 6. El árbol se puede exportar a \LaTeX{} (TikZ) o a texto.
\end{itemize}

\end{document}
